\section{Evidence architecture and reproducibility}\label{app:evidence}

This edition preserves five external analytical reports dated 5 October 2026: the large-\(L\) report R1~\cite{R1}, known-structure review R2~\cite{R2}, full two-parameter normal form R3~\cite{R3}, full-cover moduli report R4~\cite{R4}, and elliptic closure report R5~\cite{R5}. Their original files remain unchanged. The snapshots retain their original wording, including statements subsequently refined by another report. They are evidence of the reconstruction sequence, not substitutes for the final arguments here. In particular the earlier physical-analogue review's open follow-up is resolved only to the bounded negative extent stated in Section~\ref{sec:physics}.

The paper is built directly from the editable LaTeX source, with vector mathematical figures generated by a deterministic script. The supplementary directory contains the bibliography in BibTeX and rendered-reference forms, a historical source index, the complete machine-readable claim ledger, exact source hashes, verification code and results, and a build record. The repository baseline is recorded before authoring. Files in the scientific kernel, UI, and historical corpus are hashed for before/after comparison. No commit or publication operation is part of this edition.

\subsection{Independent computational checks}
The verification script evaluates 77 exact symbolic identities with SymPy. They include reciprocity, the historical \(C\)-extension, level-set differences, resultants, derivatives, parameter transformations, the coherent normal form, the marked-fiber identity, intermediate-chart actions, scaling and quotient-chart limits, elliptic transformations, discriminants, the two-isogeny, enhanced-locus factorization, and exact historical \(j\)-values. Rational identities are reduced algebraically; numerical tolerance is not used to declare them true.

The branch-cycle implementation explicitly generates the permutation group from Eq.~\eqref{eq:cycles}, confirming transitivity, order sixteen, and identity product. This checks the finite presentation used in the proof, rather than searching for a group from approximate data.

At 100-decimal-digit working precision, twelve fixed parameter pairs and four fixed nonreal source points per pair give 48 normal-form comparisons. They include positive and negative \(\Delta\), negative \(A\), \(A=0\), \(A=-1\), a constant-map case, both signs of \(B\), and a case close to cancellation. The maximum normalized discrepancy is \(2.315\times10^{-100}\). These points are deliberately off the denominator zeros and do not assert numerical stability at a pole. The four historical examples are evaluated at the same precision; their three formulas for \(j\) agree to working precision. The displayed decimals in the body are rounded, whereas the supplementary JSON retains substantially more digits.

For the intermediate limit we use the six fixed points \(s=\pm1/2,\pm\sqrt2,\pm3\). The following table reports absolute profile errors, together with relative errors for two matching paths \(\delta=L^{-1/4}\) (bulk side) and \(\delta=L^{-3/4}\) (lock side). The exact uniform estimate in Theorem~\ref{thm:matching}, not the table, proves matching.
\begin{table}[htbp]\centering\small
\begin{tabular}{rrrr}\toprule
\(L\) & \(\max|P_h-W|\) & Bulk relative error & Lock relative error\\\midrule
\(10^4\) & \(1.1324\times10^{-1}\) & \(2.3338\times10^{-2}\) & \(1.1803\times10^{-1}\)\\
\(10^8\) & \(1.1003\times10^{-3}\) & \(2.0303\times10^{-4}\) & \(1.0153\times10^{-2}\)\\
\(10^{12}\) & \(1.1000\times10^{-5}\) & \(2.0030\times10^{-6}\) & \(1.0015\times10^{-3}\)\\\bottomrule
\end{tabular}\caption{Bounded high-precision consistency checks of the proved scaling formulas. Relative errors compare \(f_L+1\) with \(-\delta\) and \(-2/(L\delta)\), respectively.}\label{tab:numerics}
\end{table}

\subsection{Claim ledger summary}
The full ledger assigns each claim a proof location, an analytical report, a historical witness where applicable, a literature anchor where used, and an evidence status. The following summary makes the distinction visible in the manuscript itself. A dash in the history column means the modern claim is not attributed to the historical document.

\begin{longtable}{L{13mm}L{48mm}L{39mm}L{45mm}}
\caption{Evidence and proof map. ``Exact'' refers to the proof in this manuscript; historical statements are documentary assertions about a source.}\label{tab:ledger}\\
\toprule ID & Claim & Manuscript proof & Evidence context\\\midrule\endfirsthead
\toprule ID & Claim & Manuscript proof & Evidence context\\\midrule\endhead
C01 & Minus-sign target; conflicting labels and later plus sign & \S\ref{sec:history} (documentary) & H1 pp.~13--14; H2 PDF, code and CSVs; H3 p.~3; H4 p.~1.\\
C02 & Reciprocity and lock fibers & Prop.~\ref{prop:levels} & H1, H3; R3; M1. Exact algebra.\\
C03 & Roots, cancellations, and real regimes & Prop.~\ref{prop:divisors}; \S\ref{sec:elementary} & R3; corrects H1 p.~14. Exact algebra.\\
C04 & Equal-value involution and coherent normalization & Thm.~\ref{thm:normal} & R1, R3. Exact algebra.\\
C05 & Universal labeled \(V_4\) quotient & Thm.~\ref{thm:belyi} & R1--R3; standard Belyi/orbifold species~\cite{SV,KO}.\\
C06 & Completeness of reduced \(\lambda\) & \S\ref{sec:belyi} & R3. Equivalence relation explicitly fixed.\\
C07 & Real contour differs from complex equivalence & \S\ref{sec:real} & R3, R4. Exact domain restrictions.\\
C08 & Large-\(L\) charts and matching & Thm.~\ref{thm:matching}; \S\ref{sec:scaling} & R1. Proof plus bounded numeric checks.\\
C09 & Distinct scaled origins of the involutions & Eqs.~\eqref{eq:Rh}--\eqref{eq:shoulders} & R1 follow-ups. Exact finite-scale identities.\\
C10 & Degree-eight branch data and lift exceptions & Thm.~\ref{thm:xcover}; \S\ref{sec:xcover} & R3, R4. Valuation and ramification proofs.\\
C11 & Selected coarse moduli \(M_{0,5}\), four presentations & Thm.~\ref{thm:moduli} & R4; marked-target convention~\cite{Mnev}.\\
C12 & Cover type and stack qualification & \S\ref{sec:moduli}; \S\ref{sec:galois} & R4. Pair types and intrinsic intermediate.\\
C13 & Normal closure, genus and monodromy & Thm.~\ref{thm:galois}; Eq.~\eqref{eq:cycles} & R4, R5. Exact function-field and group arguments.\\
C14 & Stable-target and admissible-cover boundaries & \S\ref{sec:boundary} & R4; standard framework~\cite{Braungardt,CMR}.\\
C15 & Edwards, Montgomery, Jacobi, Legendre models & \S\ref{sec:edwards} & R5; named models~\cite{Edwards,Morain,EFD}.\\
C16 & Elliptic group action and two-torsion quotients & \S\ref{sec:group} & R5; addition law~\cite{EFD}.\\
C17 & \(\nu\) as \(X_0(4)\) coordinate; two-isogeny & Thm.~\ref{thm:X04}; Eq.~\eqref{eq:isogenous} & R5; modular interpretation~\cite{Sutherland,Morain}.\\
C18 & Lifted \(\lambda\)-orbit and divisor class & \S\ref{sec:divisor} & R5. Exact quotient/divisor argument.\\
C19 & Sphere, elliptic torus, and revolution differ & \S\ref{sec:topology} & R4, R5; historical H1/H4 interpretation delimited.\\
C20 & Historical invariants and CM status & \S\ref{sec:examples} & R3, R5; exact CM criteria~\cite{DLR}. \((\pi,e)\) unresolved.\\
C21 & Reduced spectral comparison and bounded negative dictionary & \S\ref{sec:physics} & R2, R3; chain derivation~\cite{Tong}. No full observable established.\\
C22 & No priority claim for the assembled reconstruction & \S\ref{sec:intro}; \S\ref{sec:lessons} & Dedicated prior-art audit pending.\\
\bottomrule
\end{longtable}

\section{Limits of this edition and exact claims for a subsequent audit}\label{app:audit}

The documentary record leaves the intended meaning of some historical labels unresolved; both formulas are retained rather than harmonized. Printed dates have been checked against the archived PDFs, but live Zenodo deposit metadata was not verified in this edition. The historical stability and curvature experiments are indexed, not independently rerun. There is no broad novelty search, no arithmetic extension of the unresolved \((\pi,e)\) CM question, and no physical model added to the uniform chain.

The next prior-art audit should test the following precisely delimited statements, rather than asking whether reciprocity or elliptic curves in general are new:
\begin{enumerate}
\item The exact embedding of this reversed-quadratic family into \(F=(r+k)/(r-k)\), with \(u=\rho(y-1)/(y+1)\), its equal-value involution, and the interpretation of \(\lambda\) under the specified labeled equivalence.
\item The full squaring tower \(Q_t\), its \(R\)-paired Hurwitz type, the coarse \(M_{0,5}\) classification, and the four finite coefficient presentations.
\item The explicit genus-one normal closure \(tx^2z^2-x^2-z^2+t=0\), its degree-sixteen labeled quotient, and the interpretation of the two sphere-level \(V_4\) actions through two-isogenous Kummer quotients.
\item The particular parameter dictionary \(t=B/(A+1)\), \(\nu=(1-t^2)^{-1}\), the selected cyclic order-four subgroup, the two-isogenous Legendre parameter, and the lifted regular-fiber divisor \(D_\lambda\).
\item The diagonal multiscale reconstruction, especially the rigorously matched \(W(s)=-s-2/s\) profile and the exact finite-\(L\) origins of its limiting involutions.
\item The historical provenance reconciliation and correction of pole/lock claims, which is a documentary contribution distinct from mathematical priority.
\end{enumerate}
The audit must compare the same decorations and equivalences. A known elliptic equation can settle a model identification while leaving a particular coordinate tower to be traced. Conversely, a change of variables may show that the whole formulation is standard. The outcome of that audit remains open.
